% ============================================================
% Kapitel 5: Sicherheits- und Fehlerbetrachtung
% ============================================================
\chapter{Sicherheits- und Fehlerbetrachtung}

\section{Fehlerquellen im Umgang mit LLMs}
Als eine der gefährlichsten Fehlerbilder ist vorallem das Thema prompt injecting zu nennen. Dadurch kann durch nicht befugtes Personal fahrangaben getätigt werden und dadurch dass es keine totmantaster wie bei herkömmlichen robotikliefereant taktilen steuerung mittels touchpad gibt ist automatikbetrieb schwierig. Denn die seit 2025 überarbeitete Norm 102118-2 sieht für leichtbauroboter mit wenig Gefahrenpotential. als die 102118-1 für ale anderen roboter

\section{Sicherheits- und Datenschutzanforderungen}


\section{Zusammenfassung}

Abschließend fasst das letzte Kapitel die zentralen Ergebnisse dieser Arbeit zusammen und gibt einen Ausblick auf weiterführende Forschungsfelder.
